% Plantilla para documentos LaTeX para enunciados
% Por Pedro Pablo Aste Kompen - ppaste@uc.cl
% Licencia Creative Commons BY-NC-SA 3.0
% http://creativecommons.org/licenses/by-nc-sa/3.0/

\documentclass[12pt]{article}

% Margen de 1 pulgada por lado
\usepackage{fullpage}
% Incluye gráficas
\usepackage{graphicx}
% Packages para matemáticas, por la American Mathematical Society
\usepackage{amssymb}
\usepackage{amsmath}
% Desactivar hyphenation
\usepackage[none]{hyphenat}
% Saltar entre párrafos - sin sangrías
\usepackage{parskip}
% Español y UTF-8
\usepackage[spanish]{babel}
\usepackage[utf8]{inputenc}
% Links en el documento
\usepackage[hidelinks]{hyperref}
\usepackage[shortlabels]{enumitem}

\usepackage{algorithmic}
\usepackage{algorithm}

\newcommand{\bcap}{\ \hat{\cap} \ }
\newcommand{\bcup}{\ \hat{\cup} \ }
\newcommand{\bagP}{\hat{\mathcal{P}}}
\newcommand{\la}{\leftarrow}


\newcommand{\N}{\mathbb{N}}
\newcommand{\Z}{\mathbb{Z}}
\newcommand{\Exp}[1]{\mathcal{E}_{#1}}
\newcommand{\List}[1]{\mathcal{L}_{#1}}
\newcommand{\EN}{\Exp{\N}}
\newcommand{\LN}{\List{\N}}
\newcommand{\lrinline}{\text{\textbf{ --- }}}
\newcommand{\lrmatrix}{\text{\textbf{---}}}

\newcommand{\comment}[1]{}
\newcommand{\lb}{\\~\\}
\newcommand{\eop}{_{\square}}
\newcommand{\hsig}{\hat{\sigma}}
\newcommand{\ra}{\rightarrow}
\newcommand{\lra}{\leftrightarrow}
\newcommand{\com}{\textquotedbl}
\newcommand{\htau}{\hat{\tau}}
\newcommand{\twopartdef}[6]
{
	\left\{
		\begin{array}{ll}
			#1 &  #2 \\
			#3 &  #4 \\
			#5 &  #6
		\end{array}
	\right.
}

\newcommand{\mtwopartdef}[8]
{
	\left\{
		\begin{array}{ll}
			#1 &  #2 \\
			#3 &  #4 \\
			#5 &  #6 \\
			#7 &  #8
		\end{array}
	\right.
}

% Comentar aligns
\newcommand{\step}[1]{%
  \text{\phantom{(#1)}} \tag{#1}
}

\usepackage[sol]{optional}

\begin{document}
\thispagestyle{empty}
% Membrete
% PUC-ING-DCC-IIC1103
\begin{minipage}{2.3cm}
\includegraphics[width=2cm]{logo.pdf}
\vspace{0.5cm} % Altura de la corona del logo, así el texto queda alineado verticalmente con el círculo del logo.
\end{minipage}
\begin{minipage}{\linewidth}
\textsc{\raggedright \footnotesize
Pontificia Universidad Católica de Chile \\
Departamento de Ciencia de la Computación \\
IIC1253 - Matemáticas Discretas \\}
\end{minipage}


% Titulo
\begin{center}
\vspace{0.5cm}
{\huge\bf Tarea 5}\\
\vspace{0.2cm}
2 de noviembre de 2020\\
\vspace{0.2cm}
\footnotesize{2º semestre 2020 - Profesores G. Diéguez - F. Suárez}
\vspace{0.2cm}
\rule{\textwidth}{0.05mm}
\end{center}

% Cuerpo del documento ejemplo

\section*{Requisitos}
\begin{itemize}

  \item La tarea es individual. Los casos de copia serán sancionados con la
   reprobación del curso con nota $1.1$.

  \item \textbf{Entrega}: Hasta las 23:59:59 del 2 de noviembre a través del buzón habilitado en el sitio del curso (Canvas).

  \begin{itemize}

    \item Esta tarea debe ser hecha completamente en \LaTeX. Tareas hechas a mano o
    en otro procesador de texto \textbf{no serán corregidas}.

    \item Debe usar el template \LaTeX\ publicado en la página del curso.

    \item Cada problema debe entregarse en un archivo independiente de las demas preguntas.

    \item Los archivos que debe entregar son un archivo \texttt{PDF} por cada pregunta
    con su solución con nombre \texttt{numalumno-P1.pdf} y \texttt{numalumno-P2.pdf} ,
    junto con un \texttt{zip} con nombre \texttt{numalumno.zip}, conteniendo los archivos \texttt{numalumno-P1.tex} y \texttt{numalumno-P2.tex} que compilan su tarea.
    Si su código hace referencia a otros archivos, debe incluirlos también.

  \end{itemize}

  \item El no cumplimiento de alguna de las reglas se penalizará con un descuento de 0.5 en la nota final (acumulables).

  \item No se aceptarán tareas atrasadas.

  \item Si tiene alguna duda, el foro de Canvas es el lugar oficial para realizarla.
\end{itemize}


\newpage
\section*{Problemas}

    \subsection*{Problema 1 - Correctitud y Complejidad}
      El algoritmo MergeSort tiene una función de recurrencia dada por 

\[
   F(n) = 
 \begin{cases}
     1,              & \text{if } n < 2\\
    \lceil F(\frac{n}{2}) \rceil + \lfloor F(\frac{n}{2}) \rfloor + n, & \text{if } n \geq 2
 
\end{cases}
\]

Si reemplazamos $n$ = $2^k$ para $n$ $\geq$ $2$:


\begin{align}
    F(2^k) &= F(2^{k-1}) + F(2^{k-1}) + 2^k\\
           &= 2F(2^{k-1}) + 2^k\\
           &= 2 \cdot 2F(2^{k-2}) + 2 \cdot 2^k\\
           &= 2^2 \cdot F(2^{k-2}) + 2 \cdot 2^k
\end{align}

si lo generalizamos, para un k cualquiera:

\begin{align}
    F(2^k) &= 2^k \cdot F(1) + k \cdot 2^k
\end{align}

lo cual nos sirve pues no es una función recursiva, y podemos conocer su complejidad. Si volvemos a la variable n, como $n$ = $2^k$, $k$ = $\log_2 {n}$ y por lo tanto

\begin{align}
    F(n) = n \cdot F(1) + \log_2 n \cdot n
\end{align}

y como vimos en clases, $n\log_2 n$ prevalece sobre $n$, por lo que mergeSort tiene una complejidad de $\mathcal{O}(n\log{}n)$
     \newpage
     \subsection*{Problema 2 - Notación asintótica}
      \begin{enumerate}
    \item Para demostrar que esta afirmación es falsa, podemos demostrar que
    
    \begin{align}
        cn\sqrt{n} < \sqrt{n!}
    \end{align}
    
    a partir de cierto $n$ y para una constante arbitraria. Si reescribimos la expresión anterior, nos queda 
    
    \begin{align}
        cn < \sqrt{(n-1)!}
    \end{align}
    
    Notemos que la siguiente propiedad se cumple para $n$ a partir de $8$:
    
    \begin{align}
        (n-1)! > n^4
    \end{align}
    
    y también que $\forall n > c$ con $n \in \mathds{N}_+$
    
    \begin{align}
        n^2 = \sqrt{n^4} > cn 
    \end{align}
    
    por lo tanto, 
    
    \begin{align}
        \sqrt{(n-1)!} > cn
    \end{align}
    
    a partir de $n = 8$ y para cualquier c (con c = 1 nos basta para la demostración). Por lo tanto, 
    
    \begin{align}
        \sqrt{n!} \notin O(n\sqrt{n})
    \end{align}
    
    \item De manera trivial sabemos que 
    
    \begin{align}
        \sqrt{n} < \frac{n}{2} \hspace{4mm} \forall n \in \mathds{N} \land n > 4
    \end{align}
    
    por lo tanto,
    
    \begin{align}
        n^{\sqrt{n}} &< n^{\frac{n}{2}} \hspace{4mm} \forall n \in \mathds{N} \land n > 4 \\
        n^{\sqrt{n}} &< cn^{\frac{n}{2}} \hspace{2mm} \forall n \in \mathds{N} \land n > 4 \land c \geq 1
    \end{align}
    
    queda demostrado que la propiedad se cumple.
\end{enumerate}

\end{document}
